\originalchapter{导体、电容与介质}\label{ch:capacitor}
\originalsection{静电平衡与屏蔽}
金属中若有持续电场, 自由电荷会继续移动. 所以静电平衡时导体材料内部$\ee=0$, 整块相连导体等势, 多余电荷分布在表面. 外表面电场没有切向分量, 否则电荷沿面移动. 跨表面扁柱给出外部法向场$E_n=\sigma/\eps$.

“材料内部”不包括空腔. 若空腔内放入电荷$q$, 在导体材料中围住空腔的高斯面通量为零, 因而内表面感应电荷总量为$-q$. 若导体总电荷原为$Q$, 外表面总量便为$Q+q$. 空腔无内部电荷且边界等势时, 静电场为零; 有内部电荷时腔内场一般非零. 这把屏蔽外部静电场与消除所有内部场区别开来. Faraday冰桶实验通过将带电物移入桶腔, 检验感应电荷、接触转移和总电荷守恒.

空腔零场可通过唯一性验证. 无电荷腔内$\nabla^2V=0$, 边界等势$V_0$. 令$w=V-V_0$, 对$\nabla\cdot(w\nabla w)=|\nabla w|^2+w\nabla^2w$积分, 边界$w=0$使边界项为零, 得$\int|\nabla w|^2\dd V_{\rm vol}=0$. 因而梯度处处为零, $V=V_0$. 此处引用散度积分等于边界通量的积分定理, 其含义与第\ref{ch:math}章的小体积通量一致. 假设腔内势光滑且边界正常, 没有从“包围电荷为零”单独推出零场.

\originalsection{电容的几何含义}
两导体带$\pm Q$, 势差大小$\Delta V$, 在线性静电模型中$Q$与$\Delta V$成正比, 电容定义为$C=Q/\Delta V$, 单位法拉(F). 它表示在给定势差下能储存多少分离电荷, 由几何和材料决定. 平行板面积$A$, 间距$d$, 忽略边缘时$E=Q/(\eps A)$, 所以$\Delta V=Ed$、$C=\eps A/d$.

球形电容器内外半径$a<b$, 从场积分得$\Delta V=Q(1/a-1/b)/(4\pi\eps)$, 所以$C=4\pi\eps ab/(b-a)$. 同轴柱电容器长$\ell$, 忽略端部时$\Delta V=Q\ln(b/a)/(2\pi\eps\ell)$, 所以$C=2\pi\eps\ell/\ln(b/a)$. 不应直接拿平行板公式套曲面, 因为曲面之间场强一般不恒定.

并联电容两端势差相同, 总电荷相加, 得$C_{\rm eq}=\sum C_i$. 串联电容的中间孤立节点起初不带净电荷, 两侧感应电荷等量相反, 所以各电容电荷大小相同, 总势差相加, 得$1/C_{\rm eq}=\sum1/C_i$. 带初始净电荷的中间节点需另外写电荷守恒, 不能无条件使用该推导.

曲面电容的共同步骤是先算场、再积分势差、最后除电荷. 球壳间$\Delta V=\int_a^bQ\dd r/(4\pi\eps r^2)=Q(1/a-1/b)/(4\pi\eps)$; 柱壳间$\Delta V=\int_a^bQ\dd s/(2\pi\eps\ell s)=Q\ln(b/a)/(2\pi\eps\ell)$. 同一端点方向和同一电荷符号须贯穿两步.

\originalsection{储能与极化}
从零电荷准静态充电, 当已有$q$时势差为$q/C$, 搬运$\dd q$所需功$\dd W=(q/C)\dd q$. 因而$U=Q^2/(2C)=C\Delta V^2/2$. 在平行板中代入$Q=\eps EA$, 得$U=\eps E^2Ad/2$, 所以真空场能密度$u_E=\eps E^2/2$. 对一般静电分布, 将$U=\tfrac12\int\rho V\dd V_{\rm vol}$用$\rho=\eps\nabla\cdot\ee$及$\ee=-\nabla V$分部积分, 在适当无穷远衰减使边界项消失时也得该能密度. 有外加电源或固定电势边界时须把电源做功另计.

介质在外场中发生分子偶极取向或电荷位移, 形成束缚电荷. 在线性、均匀、各向同性材料中, 极化强度$\vct P=\eps\chi_e\ee$, 电位移$\vct D=\eps\ee+\vct P=\varepsilon\ee$, $\varepsilon=\kappa\eps$, $\kappa=1+\chi_e$. $\vct D$的通量只计自由电荷. 这不是任意材料的普遍线性规律, 铁电、非线性或各向异性材料需更一般关系.

电容被介质完全填充时$C'=\kappa C$. 若先断开电源则$Q$固定, 场与势差减为$1/\kappa$, 储能减为$1/\kappa$. 若电源仍连接则势差固定, 自由电荷增至$\kappa Q$, 储能增至$\kappa U$, 新能量来自电源. 局部填充要判断几何: 横向并排等于并联、沿场厚度叠层等于串联, 不能只按体积平均$\kappa$.


一般场能的分部积分链为
\[
\nabla\cdot(V\ee)=\nabla V\cdot\ee+V\nabla\cdot\ee=-E^2+V\rho/\eps,
\]
因而$\tfrac12\int\rho V\dd V_{\rm vol}=\tfrac\eps2\int E^2\dd V_{\rm vol}+\tfrac\eps2\oint V\ee\cdot\dd\vct A$. 对有限的孤立电荷取$V(\infty)=0$, 表面项随边界移到无穷远消失, 才得到纯场能积分. 理想点电荷仍有自能奇点, 不能用这个积分给出有限点电荷自能.

介质还有击穿阈值: 场过强时材料失去绝缘性质, 形成电荷贯通或破坏. 增大介电常数与提高击穿强度是两种不同材料性质, 不能用较大的$\kappa$单独保证安全耐压, 也不能断言任意介质都比真空更耐压.
